\par\noindent\textbf{(1)} 平均场单自旋配分函数 \(Z_1=2\cosh[\beta(h+Jzm)]\)。自旋平均 \(m=(1/Z_1)\sum_{s=\pm1}s e^{\beta s(h+Jzm)}=\tanh[\beta(h+Jzm)]\)。右侧包含自身平均值，必须解自洽而不能将 \(m\) 视为外部给定常量。
\par\noindent\textbf{(2)} 令 \(h=0\)，小 \(m\) 展开 \(\tanh x=x-x^3/3+O(x^5)\)。自洽式变成 \(0=[\beta Jz-1]m-(\beta Jz)^3m^3/3+\cdots\)。临界温度 \(k_BT_c=Jz\)，\(T>T_c\) 仅稳定零磁化，\(T<T_c\) 出现对称非零支。
\par\noindent\textbf{(3)} 当 \(T\uparrow T_c^-\) 时 \(\beta Jz-1\approx(T_c-T)/T_c\)，非零磁化满足 \(m^2\approx3(\beta Jz-1)/(\beta Jz)^3\sim3(T_c-T)/T_c\)。因此 \(m\propto(T_c-T)^{1/2}\)，平均场序参量临界指数 \(\beta_{\mathrm{MF}}=1/2\)。
\par\noindent\textbf{(4)} 对 \(m=\tanh[\beta(h+Jzm)]\) 在 \(h=0,m=0\) 求导得 \(\chi=\partial m/\partial h=\beta(1+Jz\chi)\)。所以 \(\chi=\beta/(1-\beta Jz)=1/[k_B(T-T_c)]\)（以自旋磁化对能量场计）。\(T\downarrow T_c^+\) 时发散。
